\documentclass[12pt,DIV14,BCOR10mm,a4paper,twoside,parskip=half-,headsepline,headinclude,english,ngerman,bibliography=totocnumbered]{scrreprt}
% Grundgröße 12pt, zweiseitig
\usepackage[headsepline,automark]{scrpage2}
\KOMAoptions{headinclude}
% Seitenköpfe automatisch
% Sprachpaket für Deutsch (Umlaute, Trennung,deutsche Überschriften)
\usepackage{babel}
\usepackage{blindtext}
%Graphikeinbindung, Hyperref (alles klickbar, Bookmarks),
%Math. Symbole aus AmsTeX
\usepackage{graphicx,hyperref,amssymb,url,atbegshi,etoolbox}
\usepackage[backend=biber,natbib=true, style=numeric, citestyle=numeric, maxnames=7, minnames=1, maxcitenames=2, url=true]{biblatex}
% Umlaute und über Tastatur einzugeben
\usepackage[utf8]{inputenc}
% add table of contents and other tables automatically
\usepackage[nottoc]{tocbibind}
%\usepackage[toc,page]{appendix}

\usepackage[autostyle=true,german=quotes]{csquotes}
%% json related packages
% dont use bera, can't get it at the work PC
%\usepackage{bera} % optional: just to have a nice mono-spaced font
\usepackage{listings}
\usepackage{xcolor}
\usepackage{textcomp}

%%% Change settings for epigraphs
%%% Change chapter format, if needed.
% customize chapter format:
% \KOMAoption{headings}{twolinechapter}
% \renewcommand*\chapterformat{\thechapter\autodot}

% customize dictum format:
\setkomafont{dictumtext}{\itshape\small}
\setkomafont{dictumauthor}{\normalfont}
\renewcommand*\dictumwidth{\linewidth}
\renewcommand*\dictumauthorformat[1]{--- #1}
\renewcommand*\dictumrule{}

\renewcommand{\lstlistingname}{Quellcode}

\makeatletter
\makeatother

% create glossaries
\usepackage[toc,section=chapter,nopostdot,acronym,style=list]{glossaries}
% style & idea copied from: https://tex.stackexchange.com/questions/8946/how-to-combine-acronym-and-glossary
\usepackage{xparse}
\DeclareDocumentCommand{\newdualentry}{ O{} O{} m m m m } {

  %%% The glossary entry the acronym links to
  \newglossaryentry{g-#3}{name={#5},
      description={#6},#1}

  %%% define the acronym and use the see= option
  \newglossaryentry{#3}{type=\acronymtype, name={#4}, description={#5}, first={#5 (#4)\glsadd{g-#3}}, see=[Glossareintrag:]{g-#3}}
}

\makeglossaries

%\loadglsentries[main]{acronyms.tex}

% Festlegung Kopf- und Fußzeile
\defpagestyle{meinstil}{%
{\headmark \hfill}
{\hfill \headmark}
{\hfill \headmark\hfill }
(\textwidth,.4pt)
}{%
(\textwidth,.4pt)
{\pagemark\hfill Dennis Grabowski, Julius Zint, Philip Matesanz, Torben Voltmer}
{Version 1.0 vom \today \hfill \pagemark}
{Version 1.0 vom \today\hfill\pagemark}
}
\pagestyle{meinstil}

\raggedbottom
\renewcommand{\topfraction}{1}
\renewcommand{\bottomfraction}{1}

\bibliography{literature.bib}
\graphicspath{{./images/}}

%%%%%%%%%%%%%%%%%%%%%%%%%%%%%%%%%%%%%%%%%%%%%%%%%%%%%%%%%%%%%%%%%%%%%%%%%%
\begin{document}    % hier gehts los
  \thispagestyle{empty} % Titelseite
\includegraphics[width=0.2\textwidth]{Wortmarke_WI_schwarz}

   {  ~ \sffamily
  \vfill
  {\Huge\bfseries Bericht: Maßnahmen zur Qualitätssicherung}
  \bigskip

  {\Large
  Dennis Grabowski, Julius Zint, Philip Matesanz, Torben Voltmer \\[2ex]
  Masterprojekt \enquote{Entwicklung und Analyse einer sicheren Web-Anwendung} \\
  Wintersemester 18/19
 \\[5ex]
   \today }
}
 \vfill

  ~ \hfill
  \includegraphics[height=0.3\paperheight]{H_WI_Pantone1665}

\vspace*{-3cm}

\tableofcontents  % Inhaltsverzeichnis

\chapter{Benutzte Methoden}
\section{Vier-Augen-Prinzip}
\section{Tests}

\subsection{Funktionale Tests}

Classcoverage, Branchcoverage, Linecoverage, Methodcoverage
End-to-End? Unittests? Wird DB getestet?
Welche Funktionen wurden nicht getestet?

\subsection{Tests zur Pruefung von sicherheitsrelevanten Funktionen}

Sessionmanagement?
Ist unser Cookie sicher und wirklich zufaellig?
Brute-force Schutz (Stichwort: "Firewall")
Authorization? Koennen Endpoints nur von den richtigen Privilegien angesurft werden?
Wird CSRF richtig genutzt?
Salted bcrypt richtig? Nutzen wir ausreichend Iterationen oder koennen wir mit einem Test das irgendwie knacken?
Kommt das Captcha zur richtigen Zeit? Ist das Captcha austricksbar?
Haben wir Quality rules fuer passwoerter?
Koennen wir kaputte Daten in die Forms eintragen?
Koennen wir irgendwelche Requests forgen?
Koennen SQL-Injections ausgefuehrt werden?
Setzen wir alle richtigen Security-Header?

\subsection{Penetrationstests}

Welche Penetrationstests wurden durchgefuehrt?
Gibt es ggf. automatisierbare Skripts?
Wenn ja, fuer welche Angriffe?

\section{Reviews}

Wer hat wann welchen Code reviewed?
Zeitliche Schaetzung?

\section{Static Analysis}

Welche Tools haben wir benutzt?
Wie haben wir sie eingeschaetzt?
Warum haben wir sie verwendet?
Wie haben wir sie verwendet?

\section{Quelltextinspektion}

Haben wir ggf. irgendwelche Checklisten nutzen koennen?
Siehe

\chapter{Gefundene Maengel}

\printbibliography

% Can be used to add a list of acronyms with their description
%\glsaddall
%\deftranslation{to=German}{Acronyms}{Abkürzungsverzeichnis}
%\deftranslation{to=German}{Glossary}{Glossar}
\printacronyms[title=Abkürzungsverzeichnis,toctitle=Abkürzungsverzeichnis]
\printglossary[type=main]

%\addcontentsline{toc}{chapter}{\listfigurename}
\listoffigures      % Abbildungsverzeichnis

%s\addcontentsline{toc}{chapter}{\listtablename}
% \listoftables       % Tabellenverzeichnis

\end{document}
